% ============================================================
%  Divisibilidade — Apostila Premium | PratiqueJa
%  Engine: XeLaTeX  |  Font: Inter via fontspec
% ============================================================
\documentclass[12pt]{article}

\usepackage[brazil]{babel}

\usepackage{fontspec}
\setmainfont[
  BoldFont       = Inter-Bold,
  ItalicFont     = Inter-Italic,
  BoldItalicFont = Inter-BoldItalic,
]{Inter-Regular}
\usepackage{microtype}

\usepackage[
  a4paper,
  top=1.5cm,
  bottom=2.8cm,
  left=1.5cm,
  right=1.5cm,
]{geometry}
\setlength{\footskip}{1.3cm}

\usepackage[dvipsnames,svgnames,table]{xcolor}

\definecolor{primary}{HTML}{2563EB}
\definecolor{primarydark}{HTML}{1E40AF}
\definecolor{primarylight}{HTML}{DBEAFE}
\definecolor{heroback}{HTML}{1D4ED8}
\definecolor{lightblue}{HTML}{60A5FA}
\definecolor{navybrand}{HTML}{051C4E}
\definecolor{softbg}{HTML}{F0F5FF}
\definecolor{chipbg}{HTML}{EFF6FF}
\definecolor{chipborder}{HTML}{BFDBFE}
\definecolor{bodytext}{HTML}{374151}
\definecolor{titletext}{HTML}{07142F}
\definecolor{mutedtext}{HTML}{9CA3AF}
\definecolor{sepcolor}{HTML}{F3F4F6}
\definecolor{exborder}{HTML}{E5E7EB}
\definecolor{exbg}{HTML}{F8FAFC}
\definecolor{badgepurple}{HTML}{7C3AED}
\definecolor{badgegreen}{HTML}{059669}

\colorlet{iris}{primary}
\colorlet{bluebell}{lightblue}

\usepackage{amsmath,amssymb}
\usepackage{multicol}
\usepackage{setspace}
\usepackage{parskip}
\usepackage{graphicx}
\usepackage{array}

\usepackage[most]{tcolorbox}
\tcbuselibrary{skins, breakable}

\usepackage{fancyhdr}
\pagestyle{fancy}
\fancyhf{}
\renewcommand{\headrulewidth}{0pt}
\renewcommand{\footrulewidth}{0.4pt}
\fancyfoot[L]{\small\color{mutedtext}© Pratique Já}
\fancyfoot[C]{\small\color{primary}\textbf{pratiqueja.com}}
\fancyfoot[R]{\small\color{mutedtext}\thepage}

% ============================================================
%  COMPONENTES
% ============================================================

\newcommand{\brandlogo}{%
  {\fontsize{20}{24}\selectfont\bfseries
    \textcolor{navybrand}{Pratique}\textcolor{primary}{Já}}%
}

% ── Hero ─────────────────────────────────────────────────────
\newcommand{\heroheader}[3]{%
  \begin{tcolorbox}[
    enhanced, arc=0pt, boxrule=0pt,
    colback=white, colframe=white,
    left=18pt, right=18pt, top=12pt, bottom=12pt,
    borderline north={2pt}{0pt}{primary},
    borderline south={3pt}{0pt}{primary},
  ]
    \begin{minipage}[c]{0.24\linewidth}
      \centering
      \brandlogo\par
      \vspace{4pt}
      {\color{mutedtext}\fontsize{7}{8.5}\selectfont\bfseries\MakeUppercase{#3}}
    \end{minipage}%
    \hfill
    \begin{minipage}[c]{0.72\linewidth}
      \centering
      {\fontsize{26}{32}\selectfont\bfseries\color{heroback}#1}\par
      \vspace{4pt}
      {\color{primary}\rule{0.44\linewidth}{1.2pt}}\par
      \vspace{4pt}
      {\color{bodytext}\normalsize\itshape #2}
    \end{minipage}
  \end{tcolorbox}%
  \vspace{6pt}%
}

% ── Badge inline ──────────────────────────────────────────────
\newcommand{\badge}[2]{%
  \tcbox[
    on line, enhanced, arc=7pt, boxrule=0pt,
    colback=#2, colframe=#2,
    left=5pt, right=5pt, top=3pt, bottom=3pt,
    nobeforeafter,
  ]{\small\bfseries\color{white}#1}%
}

% ── Cabeçalho de seção — título e subtítulo na mesma linha ───
\newcommand{\TW}[4]{%
  \vspace{7pt}%
  \noindent{\color{#2}\rule{\linewidth}{2.5pt}}\par\nopagebreak
  \vspace{3pt}%
  \noindent\badge{#1}{#2}\hspace{6pt}%
  {\color{titletext}\bfseries\normalsize #3}%
  {\color{mutedtext}\ ·\ \small\itshape #4}\par\nopagebreak
  \vspace{3pt}%
  {\color{sepcolor}\rule{\linewidth}{1pt}}%
  \vspace{4pt}%
}

% ── Caixa de regra — .tr-rule ────────────────────────────────
\newcommand{\TM}[1]{%
  \begin{tcolorbox}[
    enhanced, arc=0pt, boxrule=0pt,
    colback=softbg, colframe=softbg,
    left=10pt, right=8pt, top=7pt, bottom=7pt,
    borderline west={4pt}{0pt}{primary},
  ]
    \color{primarydark}\small #1
  \end{tcolorbox}%
  \vspace{2pt}%
}

% ── Caixa de exemplos agrupados ───────────────────────────────
% Todos os exemplos de uma seção em uma única caixa compacta.
% Use \\ para separar linhas internamente.
\newcommand{\exemp}[1]{%
  \vspace{2pt}%
  \begin{tcolorbox}[
    arc=7pt, boxrule=0.6pt,
    colback=exbg, colframe=exborder,
    left=9pt, right=9pt, top=5pt, bottom=5pt,
  ]
    \setlength{\parskip}{2pt}%
    \small\color{bodytext}#1
  \end{tcolorbox}%
  \vspace{1pt}%
}

% ── Rótulo h4 compacto ────────────────────────────────────────
\newcommand{\SH}[1]{%
  \vspace{4pt}%
  \noindent{\color{mutedtext}\fontsize{7.5}{9}\selectfont
            \bfseries\MakeUppercase{#1}}\par
  \vspace{1pt}%
}

% ── Caixa de fórmula ─────────────────────────────────────────
\newcommand{\formula}[1]{%
  \vspace{3pt}%
  \begin{tcolorbox}[
    enhanced, arc=0pt, boxrule=0pt,
    colback=primarylight, colframe=primarylight,
    left=10pt, right=10pt, top=7pt, bottom=7pt,
    borderline west={4pt}{0pt}{primary},
  ]
    \centering\color{primarydark}\normalsize\bfseries #1
  \end{tcolorbox}%
  \vspace{3pt}%
}

% ── Chips de divisor ──────────────────────────────────────────
\newcommand{\divisorchip}[1]{%
  \tcbox[
    on line, enhanced, arc=6pt, boxrule=0.6pt,
    colback=chipbg, colframe=chipborder,
    left=5pt, right=5pt, top=2pt, bottom=2pt,
    nobeforeafter,
  ]{\small\bfseries\color{primarydark}#1}\hspace{2pt}%
}

% ============================================================
\begin{document}
\pagenumbering{arabic}
\setstretch{1.25}
\color{bodytext}

\heroheader{Divisibilidade}%
           {Regras, Divisores e Aplicações}%
           {Matemática · Intermediário}

\setlength{\columnsep}{20pt}
\setlength{\columnseprule}{0.5pt}
\def\columnseprulecolor{\color{chipborder}}

\begin{multicols}{2}

%─────────────────────────────────────────────────────────────
\TW{Def}{badgepurple}{Definição}{O que é divisibilidade?}

\TM{Um número é divisível por outro quando o \textbf{resto da divisão
é zero}. Dizemos que $a$ é divisível por $b$ quando existe $q \in
\mathbb{Z}$ tal que $a = b \cdot q$. Chamamos $b$ de \textbf{divisor}
de $a$ e $a$ de \textbf{múltiplo} de $b$.}

%─────────────────────────────────────────────────────────────
\TW{2}{primary}{Divisibilidade por 2}{Números pares}

\TM{O \textbf{último dígito} deve ser \textbf{0, 2, 4, 6 ou 8}.}

\exemp{%
  $13\textcolor{primary}{8}$ → ult.\ dígito 8 → divisível ✓\\
  $47\textcolor{primary}{5}$ → ult.\ dígito 5 → não divisível ✗%
}

%─────────────────────────────────────────────────────────────
\TW{3}{primary}{Divisibilidade por 3}{Soma dos algarismos}

\TM{A \textbf{soma dos algarismos} deve ser múltipla de \textbf{3}.}

\exemp{%
  $132 \Rightarrow 1{+}3{+}2=\textcolor{primary}{6}$ → mult.\ de 3 ✓\\
  $531 \Rightarrow 5{+}3{+}1=\textcolor{primary}{9}$ → mult.\ de 3 ✓\\
  $124 \Rightarrow 1{+}2{+}4=\textcolor{primary}{7}$ → não divisível ✗%
}

%─────────────────────────────────────────────────────────────
\TW{4}{primary}{Divisibilidade por 4}{Dois últimos dígitos}

\TM{Os \textbf{2 últimos dígitos} devem formar número divisível por 4.}

\exemp{%
  $5\textcolor{primary}{24}$ → 24 é múltiplo de 4 ✓\\
  $13\textcolor{primary}{36}$ → 36 é múltiplo de 4 ✓\\
  $5\textcolor{primary}{14}$ → 14 não é múltiplo de 4 ✗%
}

%─────────────────────────────────────────────────────────────
\TW{5}{primary}{Divisibilidade por 5}{Último dígito 0 ou 5}

\TM{O \textbf{último dígito} deve ser \textbf{0} ou \textbf{5}.}

\exemp{%
  $32\textcolor{primary}{5}$ → termina em 5 ✓\quad
  $14\textcolor{primary}{0}$ → termina em 0 ✓\\
  $78\textcolor{primary}{3}$ → termina em 3 → não divisível ✗%
}

%─────────────────────────────────────────────────────────────
\TW{6}{primary}{Divisibilidade por 6}{Divisível por 2 e por 3}

\TM{Deve ser \textbf{par} e ter soma dos algarismos \textbf{múltipla
de 3} simultaneamente.}

\exemp{%
  $654$: par e $6{+}5{+}4=\textcolor{primary}{15}$ (mult.\ 3) ✓\\
  $312$: par e $3{+}1{+}2=\textcolor{primary}{6}$ (mult.\ 3) ✓\\
  $125$: ímpar → não divisível por 6 ✗%
}

%─────────────────────────────────────────────────────────────
\TW{7}{primary}{Divisibilidade por 7}{Regra do duplo}

\TM{Multiplique o \textbf{último dígito por 2} e subtraia do número
sem esse dígito. Repita até concluir.}

\exemp{%
  $16\textcolor{primary}{1}
  \Rightarrow 1{\cdot}2=2
  \Rightarrow 16{-}2=\textcolor{primary}{14}$ ✓\\
  $109\textcolor{primary}{9}
  \Rightarrow 9{\cdot}2=18
  \Rightarrow 109{-}18=91
  \Rightarrow 9{-}2=\textcolor{primary}{7}$ ✓%
}

%─────────────────────────────────────────────────────────────
\TW{8}{primary}{Divisibilidade por 8}{Três últimos dígitos}

\TM{Os \textbf{3 últimos dígitos} devem formar número divisível por 8.}

\exemp{%
  $7\textcolor{primary}{128}$ → 128 é múltiplo de 8 ✓\\
  $5\textcolor{primary}{256}$ → 256 é múltiplo de 8 ✓\\
  $3\textcolor{primary}{100}$ → 100 não é múltiplo de 8 ✗%
}

%─────────────────────────────────────────────────────────────
\TW{9}{primary}{Divisibilidade por 9}{Soma dos algarismos}

\TM{A \textbf{soma dos algarismos} deve ser múltipla de \textbf{9}.}

\exemp{%
  $432 \Rightarrow 4{+}3{+}2=\textcolor{primary}{9}$ ✓\\
  $6561 \Rightarrow 6{+}5{+}6{+}1=\textcolor{primary}{18}$ ✓\\
  $124 \Rightarrow 1{+}2{+}4=\textcolor{primary}{7}$ → não divisível ✗%
}

%─────────────────────────────────────────────────────────────
\TW{10}{primary}{Divisibilidade por 10}{Último dígito 0}

\TM{O \textbf{último dígito deve ser 0}. Como $10=2{\times}5$,
equivale a ser divisível por 2 e por 5 ao mesmo tempo.}

\exemp{%
  $15\textcolor{primary}{0}$ → termina em 0 ✓\quad
  $487\textcolor{primary}{0}$ → termina em 0 ✓\\
  $32\textcolor{primary}{5}$ → termina em 5 → div.\ por 5, não por 10 ✗%
}

%─────────────────────────────────────────────────────────────
\TW{11}{badgegreen}{Divisibilidade por 11}{Soma alternada dos algarismos}

\TM{A \textbf{soma alternada} dos algarismos (alternando $+$ e $-$
da esquerda para a direita) deve ser múltipla de 11 ou zero.}

\exemp{%
  $253 \Rightarrow {+}2{-}5{+}3 = \textcolor{primary}{0}$
  → mult.\ de 11 ✓\\
  $1452 \Rightarrow {+}1{-}4{+}5{-}2 = \textcolor{primary}{0}$ ✓\\
  $123 \Rightarrow {+}1{-}2{+}3 = \textcolor{primary}{2}$
  → não divisível ✗%
}

%─────────────────────────────────────────────────────────────
\TW{Div}{badgegreen}{Divisores de um Número}{Como encontrar todos os divisores naturais}

Para cada fator primo da decomposição, multiplique \textbf{todos os
divisores já encontrados} por esse fator e acrescente os novos.

\SH{Exemplo — divisores de 60}

$60 = 2^2 \cdot 3^1 \cdot 5^1$

\vspace{3pt}
\noindent
$\begin{array}[t]{r|l|l}
  & & \textcolor{primary}{1} \\
60 & \textcolor{primary}{2} & \textcolor{primary}{2} \\
30 & \textcolor{primary}{2} & \textcolor{primary}{4} \\
15 & \textcolor{primary}{3} & \textcolor{primary}{3,\ 6,\ 12} \\
 5 & \textcolor{primary}{5} & \textcolor{primary}{5,\ 10,\ 20,\ 15,\ 30,\ 60} \\
\hline
 1 & & \\
\end{array}$

\SH{Todos os divisores de 60}

\noindent
\divisorchip{1}\divisorchip{2}\divisorchip{3}\divisorchip{4}%
\divisorchip{5}\divisorchip{6}\divisorchip{10}\divisorchip{12}%
\divisorchip{15}\divisorchip{20}\divisorchip{30}\divisorchip{60}

\SH{Quantidade de divisores}

Se $n = p_1^{a_1} \cdot p_2^{a_2} \cdots$, o total de divisores é:

\formula{$d(n) = (a_1+1)\cdot(a_2+1)\cdots$}

\exemp{%
  $60 = 2^2{\cdot}3^1{\cdot}5^1
  \Rightarrow d(60)=(3)(2)(2)=\textcolor{primary}{12}$ divisores\\
  $36 = 2^2{\cdot}3^2
  \Rightarrow d(36)=(3)(3)=\textcolor{primary}{9}$ divisores%
}

\end{multicols}

\end{document}
